\diarytitle{0047}{定数係数線形微分方程式の初期値問題（ラプラス変換、非同次項 e^{-t}）}

関数 $f(t)$（$t \ge 0$）のラプラス変換を $\mathcal{L}[f(t)](s) = F(s) = \displaystyle\int_{0}^{\infty} e^{-st} f(t)\,dt$ と定義する。初期条件を代入した上で方程式全体をラプラス変換すると、$\mathcal{L}[y'(t)](s) = sY(s) - y(0)$、$\mathcal{L}[y''(t)](s) = s^{2}Y(s) - s\,y(0) - y'(0)$（$Y(s) = \mathcal{L}[y(t)](s)$）により $Y(s)$ についての代数方程式が得られる。これを $Y(s)$ について解き、部分分数分解した上で基本変換表を逆に用いて逆ラプラス変換すれば $y(t)$ が求まる。

両辺をラプラス変換する。$y(0)=y'(0)=0$ より
\[
s^{2}Y(s) + 3sY(s) + 2Y(s) = \frac{1}{s+1}
\]
左辺をまとめて
\[
(s^{2}+3s+2)\,Y(s) = (s+1)(s+2)\,Y(s) = \frac{1}{s+1}
\]
よって
\[
Y(s) = \frac{1}{(s+1)^{2}(s+2)}
\]
部分分数分解を
\[
Y(s) = \frac{A}{s+1} + \frac{B}{(s+1)^{2}} + \frac{C}{s+2}
\]
とおくと、$1 = A(s+1)(s+2) + B(s+2) + C(s+1)^{2}$。
$s=-1$ を代入して $B=1$、$s=-2$ を代入して $C=1$。$s^{2}$ の係数を比較して $0 = A+C$ より $A=-1$。したがって
\[
Y(s) = -\frac{1}{s+1} + \frac{1}{(s+1)^{2}} + \frac{1}{s+2}
\]
逆ラプラス変換して
\[
\boxed{\; y(t) = -e^{-t} + t\,e^{-t} + e^{-2t} \;}
\]
検算: $y(0) = -1+0+1=0$。$y'(t) = 2e^{-t} - t e^{-t} - 2e^{-2t}$ より $y'(0) = 2-0-2=0$。$y''(t) = -3e^{-t}+t e^{-t}+4e^{-2t}$ であり、
\[
y''+3y'+2y = (-3e^{-t}+te^{-t}+4e^{-2t}) + 3(2e^{-t}-te^{-t}-2e^{-2t}) + 2(-e^{-t}+te^{-t}+e^{-2t}) = e^{-t}
\]
となり、右辺 $e^{-t}$ と一致する。
